\documentclass[conference]{IEEEtran}
\usepackage[utf8]{inputenc}
\usepackage{cite}
\usepackage{amsmath}
\usepackage{hyperref}

\title{Technical Annex: A Toy Single-Area Frequency-Response\\
Simulator for AI-Assisted Grid Stability Control}

\author{\IEEEauthorblockN{Judith Drack}\\
\IEEEauthorblockA{Companion to ``Power System Stability in a Decarbonising Europe''}}

\begin{document}
\maketitle

\begin{abstract}
This technical annex documents an interactive single-area frequency-response simulator built to illustrate the application of artificial intelligence to the grid stability problem discussed in the companion paper. It describes (i)~the physical meaning of the simulator's user-facing parameters, (ii)~the swing-equation dynamics that govern the simulated frequency response, (iii)~the three controllers---no control, classical PID, and an illustrative reinforcement-learning-style policy---that the simulator compares, and (iv)~the diagnostic metrics (frequency nadir, rate of change of frequency, recovery time, and a composite stability score) together with the qualitative regimes those metrics distinguish. The annex is deliberately conservative about what the toy model can and cannot claim, and explicitly identifies the modelling simplifications that would have to be relaxed before any operational conclusion could be drawn.
\end{abstract}

\section{Scope}
The simulator is a deliberately small \emph{toy} model: a single state variable (system frequency) integrated by a forward-Euler scheme on the swing equation, with one aggregated governor, one battery / fast-frequency-response (FFR) asset, and one disturbance. It is intended as a pedagogical companion to the paper and is not a substitute for utility-grade tools such as PSS\textsuperscript{\textregistered}E, DIgSILENT PowerFactory, PowerWorld, or open-source frameworks such as PyPSA and PowerSimulationsDynamics.jl. The modelling simplifications and the phenomena they hide are listed in Section~\ref{sec:limits}.

\section{User-Facing Parameters}\label{sec:params}

\subsection{Synchronous-machine share}
The first slider, expressed as a percentage from 0\% to 100\%, sets the aggregate system inertia constant $H$ (in seconds) according to the linear mapping
\begin{equation}
H \;=\; 0.5 \;+\; 4.5\,\frac{\text{share}\,[\%]}{100}.
\end{equation}
A 100\% synchronous fleet yields $H = 5$\,s, close to the historical Continental European synchronous area value~\cite{tielens2016relevance,ulbig2014impact}. A 0\% share (a fully inverter-based system) yields a residual $H = 0.5$\,s, representing the small effective inertia contributed by wind-turbine kinetic energy when controllers extract it, and by converter-based synthetic-inertia provisions. Real aggregate inertia is not strictly linear in synchronous-machine share---it depends on which machines are dispatched and on their individual inertia constants---but the linear mapping is sufficient for the simulator's qualitative purpose.

\subsection{Battery / fast-frequency-response capacity}
The second slider sets the maximum battery dispatch $P_{\mathrm{bat,max}}$ in MW, between 0 and 600. This represents the headroom available to a fast frequency response asset---typically a grid-scale lithium-ion battery, but functionally including aggregated demand response or fast-acting synchronous condensers. FFR is procured as an ancillary service by several transmission system operators, with specified response time and energy-duration requirements~\cite{denholm2020inertia,nerc2015erstf}. The simulator models the battery as a first-order system with a 50\,ms time constant, faster than but not far from existing commercial deployments.

\subsection{Disturbance size}
The third slider sets the magnitude of a simulated generation loss in MW (100--1500), applied as a step at $t = 2$\,s. The 600\,MW default corresponds approximately to the loss of a single large thermal unit. ENTSO-E's design contingency for the Continental European synchronous area is a 3\,000\,MW reference incident---typically the loss of two large units, or a major HVDC interconnection---against which RoCoF and nadir withstand requirements are calibrated~\cite{entsoe2020inertia,entsoe2017hpopeips}. The simulator's system base is $S_{\mathrm{b}} = 5\,000$\,MVA, so the slider's range spans approximately 2\% to 30\% of base.

\section{Dynamics: the Swing Equation}\label{sec:swing}
The frequency of a single-area system evolves according to the \emph{swing equation}~\cite{kundur1994}:
\begin{equation}
2H\,\frac{d}{dt}\!\left(\frac{\Delta f}{f_0}\right) \;=\; \frac{\Delta P}{S_{\mathrm{b}}},
\label{eq:swing-pu}
\end{equation}
where $f_0 = 50$\,Hz is the nominal frequency and $\Delta f = f - f_0$ is the deviation. Isolating the time derivative,
\begin{equation}
\frac{d\Delta f}{dt} \;=\; \frac{f_0}{2 H\, S_{\mathrm{b}}}\,\Delta P,
\label{eq:swing-mw}
\end{equation}
which states that the rate of change of frequency is proportional to the net power imbalance, with a constant of proportionality set by inertia. The instantaneous post-disturbance RoCoF---$d\Delta f/dt$ evaluated at the moment of the contingency---is therefore almost entirely determined by $H$ and is largely independent of the control action that follows. A controller cannot meaningfully reduce RoCoF; it can only reduce the \emph{nadir} (the integral of RoCoF over the response window) by injecting power quickly enough.

\subsection{Load damping}
Loads consume slightly less real power as frequency falls---and slightly more as frequency rises---an effect captured by a damping coefficient $D \approx 1$\,p.u.~\cite{kundur1994}:
\begin{equation}
\Delta P_{\mathrm{damping}} \;=\; -D\,S_{\mathrm{b}}\,\frac{\Delta f}{f_0}.
\end{equation}

\subsection{Primary frequency control (governor)}
The synchronous-machine fleet supplies a primary frequency response through governor action, modelled as a first-order system tracking a steady-state target proportional to $-\Delta f / R$, with droop $R = 0.05$ (5\%) and time constant $T_{\mathrm{gov}} = 5$\,s:
\begin{align}
P_{\mathrm{gov,target}} &= -\frac{S_{\mathrm{b}}}{R}\,\frac{\Delta f}{f_0},\\
T_{\mathrm{gov}}\,\dot P_{\mathrm{gov}} &= P_{\mathrm{gov,target}} - P_{\mathrm{gov}}.
\end{align}

\subsection{Net imbalance and integration}
The total imbalance entering~\eqref{eq:swing-mw} is
\begin{equation}
\Delta P \;=\; P_{\mathrm{gov}} + P_{\mathrm{bat}} + P_{\mathrm{dist}} \;-\; D\,S_{\mathrm{b}}\,\frac{\Delta f}{f_0},
\end{equation}
with $P_{\mathrm{dist}} = -P_{\mathrm{disturbance}}$ for $t \geq t_{\mathrm{dist}} = 2$\,s and zero otherwise. The simulator integrates the resulting four-state system $(\Delta f, P_{\mathrm{gov}}, P_{\mathrm{bat}}, \int\!\Delta f\,dt)$ with a forward-Euler scheme of step $\delta t = 10$\,ms over 30\,s of simulated time. The Euler step is adequate for the qualitative purpose of the simulator but would not be acceptable for an engineering study, where a stiff implicit solver would be required.

\section{Controllers}\label{sec:ctrl}
The simulator implements three controllers that share the same battery actuator but differ in how they compute its setpoint.

\subsection{None}
$P_{\mathrm{bat}} = 0$. Only the governor reacts. This is the realistic baseline for a low-inertia grid without dedicated FFR provisions.

\subsection{Classical PID}
A proportional--integral--derivative controller on the measured frequency deviation:
\begin{equation}
P_{\mathrm{bat,target}}^{\mathrm{PID}} \;=\; -\bigl(K_P\,\Delta f + K_I\!\!\int\!\Delta f\,dt + K_D\,\dot{\Delta f}\bigr),
\label{eq:pid}
\end{equation}
with $K_P = 800$, $K_I = 200$, $K_D = 50$, clamped to $|P_{\mathrm{bat}}| \leq P_{\mathrm{bat,max}}$. The controller is purely reactive: it acts on the measured $\Delta f$ and its integral and derivative.

\subsection{Illustrative reinforcement-learning policy}\label{sec:rl}
A real reinforcement-learning agent for this task would be trained against the simulator (or a higher-fidelity environment) to maximise a reward such as
\begin{equation}
r_t \;=\; -\bigl(w_1\,|\Delta f_t| + w_2\,|\dot{\Delta f}_t| + w_3\,|P_{\mathrm{bat},t}|\bigr) \;-\; w_4\,\mathbb{1}\!\left[\text{UFLS}\right],
\label{eq:reward}
\end{equation}
using an algorithm such as deep Q-learning or proximal policy optimisation~\cite{glavic2017rl}. The L2RPN challenges, organised by the French transmission system operator RTE, are a public benchmark for this class of problem at network scale~\cite{marot2020l2rpn}.

In the annex's simulator we do \emph{not} train such an agent. Instead we implement a hand-tuned policy that illustrates the qualitative behaviour a learned policy would aim to exhibit: an aggressively-tuned PID ($K_P = 1400$, $K_I = 250$, $K_D = 180$), supplemented with a \emph{feed-forward burst} when a large RoCoF is detected within the first 1.5\,s after the disturbance:
\begin{equation}
P_{\mathrm{bat,target}}^{\mathrm{RL}} \;=\; P_{\mathrm{bat,target}}^{\mathrm{aggressive\,PID}} - \mathrm{sgn}(\dot{\Delta f})\,P_{\mathrm{bat,max}}\,\mathbb{1}\!\left[|\dot{\Delta f}| > 0.2,\; t - t_{\mathrm{dist}} < 1.5\,\text{s}\right].
\label{eq:rl}
\end{equation}
The feed-forward term uses the RoCoF as a proxy for ``a large disturbance is in progress'' and pre-commits the battery in the direction opposing the imbalance, rather than waiting for $\Delta f$ to integrate. This is exactly the kind of policy a trained agent would discover from the reward function in~\eqref{eq:reward}: when the cost of being late is large (UFLS), the optimal policy is to spend control effort early. The simulator labels this controller as \emph{illustrative} throughout the user interface to avoid overstating what is being demonstrated.

\section{Diagnostic Metrics}\label{sec:metrics}
For each run the simulator computes:

\subsection{Frequency nadir}
The lowest frequency $f_{\min}$ reached during the simulation. The first stage of under-frequency load shedding (UFLS) in the Continental European synchronous area is at 49.0\,Hz; falling below it triggers automatic disconnection of a fraction of load, which is already a localised, controlled partial blackout. A safe operating margin is $f_{\min} \gtrsim 49.5$\,Hz.

\subsection{Maximum RoCoF}
$\max_t |\dot{\Delta f}(t)|$. ENTSO-E design standards aim for RoCoF tolerances of order 1\,Hz/s; some converter-rich systems already tolerate up to 2\,Hz/s with adapted protection~\cite{entsoe2020inertia}. As noted above, RoCoF is dominated by the instantaneous inertia at the moment of disturbance.

\subsection{Recovery time}
The first instant after the disturbance at which $|\Delta f|$ stays below 0.1\,Hz for one full second continuously, indicating that primary and supplementary control have re-established a stable operating point. Recovery within 20--30\,s is typical for healthy primary frequency response.

\subsection{Composite stability score}
A linear penalty combination intended only as a coarse summary, not as an engineering metric:
\begin{equation}
\begin{aligned}
\text{score} &= 100 \\
&\quad - 60\,\max\bigl(0,\;(f_0 - f_{\min}) - 0.2\bigr)\\
&\quad - 25\,\max\bigl(0,\;\max_t|\dot{\Delta f}| - 0.5\bigr)\\
&\quad - 25\,\mathbb{1}[\text{UFLS}],
\end{aligned}
\end{equation}
clamped to $[0,\,100]$.

\section{Qualitative Regimes}\label{sec:regimes}
The simulator's headline observation is that the marginal value of a smarter controller increases sharply as inertia falls.

\subsection{High-inertia regime}
With a 80\% synchronous-machine share and a 600\,MW disturbance, all three controllers comfortably keep the system above UFLS; the differences in nadir among them are tenths of a hertz. The historical European grid operated in this regime, and classical primary control sufficed.

\subsection{Low-inertia regime}
With a 20\% synchronous-machine share and the same disturbance, the uncontrolled response triggers UFLS, the PID controller lands near the threshold, and the illustrative RL policy keeps the nadir comfortably above. This is the qualitative point the simulator is built to make: AI-assisted, anticipatory control can extend the operating envelope of a low-inertia grid by exploiting the early RoCoF signal that classical reactive controllers cannot use.

\subsection{Out-of-envelope regime}
With $H \to 0.5$\,s and a 1\,500\,MW disturbance, even the illustrative RL policy with 600\,MW of battery cannot save the system: there is insufficient fast response to absorb the imbalance before the very low inertia drives the frequency through the floor. This is an honest negative result: AI does not substitute for physical resources, only for the cleverness with which existing resources are dispatched.

\section{Modelling Limitations}\label{sec:limits}
The single-area assumption hides several phenomena that matter in real European studies, and no operational claim should be drawn from this simulator without addressing them:

\begin{itemize}
\item \textbf{Inter-area oscillations.} Multi-machine models with explicit tie-lines are required to study the slow swings between geographic regions of the European synchronous area.
\item \textbf{Voltage and reactive-power dynamics.} The model tracks only frequency. Voltage stability and reactive-power coordination---which become more difficult as IBR penetration grows---are out of scope~\cite{hatziargyriou2021revisited}.
\item \textbf{Protection and short-circuit behaviour.} Converters limit their fault current; the displacement of synchronous machines alters relay coordination. None of this is captured~\cite{entsoe2017hpopeips}.
\item \textbf{Stochastic generation and load.} The disturbance is a single deterministic step; renewable variability and load forecast errors are absent.
\item \textbf{Detailed converter control.} The grid-forming versus grid-following distinction~\cite{lasseter2020gridforming} is collapsed into a single battery actuator.
\item \textbf{Numerical method.} Forward-Euler integration with $\delta t = 10$\,ms is adequate for qualitative use but inappropriate for engineering studies, where a stiff implicit solver would be required.
\end{itemize}
Any engineering claim about real grid behaviour would require relaxing these simplifications and adopting utility-grade simulation tools, and any deployment claim about an AI controller would require training against a higher-fidelity environment and validating under the governance regime of the EU AI Act and NIS\,2 Directive discussed in the companion paper.

\section{Implementation Notes}
The simulator is a single-file HTML/JavaScript application that runs entirely in the browser, with approximately thirty lines of physics code and the remainder devoted to user interface and chart rendering. The source is available locally as \texttt{grid-stability-tool.html}.

\bibliographystyle{IEEEtran}
\bibliography{judith}

\end{document}
